\section{Energy Storage and Power Electronics}

\begin{multicols}{2}
The energy storage and low-voltage electrical distribution networks act as the circulatory system of the EV powertrain. The specified load profile—featuring a 4.1 kW continuous nominal draw and severe 25.2 kW transient peaks—demands a custom, highly optimized 72V Lithium-ion architecture complemented by a robust DC-DC auxiliary conversion stage.

\subsection{Battery Voltage and Capacity Sizing}

The fundamental power source is a high-energy-density Lithium-ion battery pack. To satisfy the 72V system requirement, the baseline cell selected is a standard 18650/21700 geometry Lithium-ion (NMC) cell featuring a nominal voltage ($V_{cell}$) of 3.7 V. 

To achieve the targeted operational system voltage, the required number of cells connected in series ($N_s$) is linearly derived:
\begin{equation} \label{eq:series_config}
\begin{split}
N_s &= \frac{V_{target}}{V_{cell}} \\
&= \frac{72\,\mathrm{V}}{3.7\,\mathrm{V}} \\
&\approx 19.45 \rightarrow 20\,\mathrm{S}
\end{split}
\end{equation}

Selecting a 20-Series (20S) configuration yields a true nominal pack voltage of $20 \times 3.7\,\mathrm{V} = 74\,\mathrm{V}$, providing an absolute peak full-charge voltage of 84.0 V ($20 \times 4.2\,\mathrm{V}$). This 20S architecture is conventionally classified as a "72V" class battery in lightweight EV designations.

The required energy capacity ($E_{capacity}$) is calculated directly against the vehicle's operational target autonomy. Assuming a mandate for exactly 45 minutes ($0.75\,\mathrm{h}$) of continuous running at the previously calculated 80 km/h cruising power load ($P_{cont\_motor} = 4.1\,\mathrm{kW}$):
\begin{equation} \label{eq:energy_capacity}
\begin{split}
E_{capacity} &= P_{cont\_motor} \times t_{run} \\
&= 4,100\,\mathrm{W} \times 0.75\,\mathrm{h} \\
&= 3,075\,\mathrm{Wh} \approx 3.1\,\mathrm{kWh}
\end{split}
\end{equation}

Converting this absolute energy requirement into charge capacity (Ampere-hours) at the nominal 74V pack potential:
\begin{equation} \label{eq:ah_capacity}
\begin{split}
Q_{pack} &= \frac{E_{capacity}}{V_{pack(nom)}} \\
&= \frac{3,075\,\mathrm{Wh}}{74\,\mathrm{V}} \\
&\approx 41.5\,\mathrm{Ah}
\end{split}
\end{equation}

If an individual selected cell possesses a capacity of 4,000 mAh (4.0 Ah), fulfilling this demand requires approximately 11 parallel strings ($11 \times 4.0\,\mathrm{Ah} = 44\,\mathrm{Ah}$). Therefore, the comprehensive macro-architecture of the battery is specified as a \textbf{20S11P} configuration, integrating exactly 220 independent Lithium-ion cells for a total energy capacity of 3.25 kWh.

\subsection{Comprehensive Battery Design}

Deploying a 20S11P block into a violent mechanical environment necessitates a holistically integrated tri-lateral design philosophy encompassing electrical safety, mechanical rigidity, and thermal resilience.

\textbf{1. Electrical Design:}
The parallel groups (P-groups) are interconnected utilizing pure nickel strip busbars (0.2 mm thickness $\times$ 10 mm width) spot-welded directly to the cell terminals to permanently mitigate vibrational contact resistance. The absolute peak inverter current draw ($I_{peak}$) when extracting 25.2 kW from the 74V pack is a massive 340 Amperes. To survive this load safely, the main electrical traction bus is governed by an automotive-grade Battery Management System (BMS). The BMS continuously measures individual cell voltages, governing active balancing during charging, and controlling the main electromechanical contactors to instantly isolate the pack upon detection of an overcurrent ($>350$A) or undervoltage ($<60$V) fault condition.

\textbf{2. Mechanical Design:}
Due to extreme inertial cornering and acceleration loads, the 220 individual cylindrical cells cannot bear structural tension. They are nested within flame-retardant Acrylonitrile Butadiene Styrene (ABS) cell-spacers. These spacers establish a precise 2 mm air gap between adjacent cylindrical walls, fundamentally preventing catastrophic localized propagation of thermal runaway. The entire cohesive ABS matrix is then isolated from chassis chassis flex via dense polyurethane foam mounting blocks and encased within an IP67-rated 6061-T6 aluminum exo-box.

\textbf{3. Thermal Design:}
High internal current densities generate severe, non-linear Joule heating. The individual cell internal resistance ($R_{cell}$) is roughly 15 m$\Omega$. The aggregate pack resistance ($R_{pack}$) for the 20S11P matrix is analytically determined as:
\begin{equation} \label{eq:pack_resistance}
\begin{split}
R_{pack} &= \frac{N_s \cdot R_{cell}}{N_p} \\
&= \frac{20 \cdot 0.015\,\Omega}{11} \\
&\approx 0.027\,\Omega
\end{split}
\end{equation}

During sustained maximum 340 A discharge events, the resultant $I^2R$ power rejected as internal heat ($P_{loss}$) is:
\begin{equation} \label{eq:thermal_loss}
P_{loss} = I_{peak}^2 R_{pack} = (340)^2 \cdot 0.027 \approx 3,121\,\mathrm{W}
\end{equation}

While 3.1 kW of localized heating is formidable, it occurs exclusively during rare transient peaks rather than continuous load. Therefore, forced-convection air cooling is strategically selected over complex liquid dielectrics to severely reduce mass and complexity. High static-pressure blowers directly channel ambient air through the explicitly designed 2 mm spacer channels, creating turbulent flow across the cell casing to strip heat away and maintain junction temperatures strictly below the 60$^\circ$C threshold.
\end{multicols}

\begin{figure}[H]
\centering
\includegraphics[width=\textwidth]{"Figures Images/figure4_efficiency.png"}
\caption{System block diagram highlighting the efficiency chain and the DC-DC buck step-down isolation.}
\label{fig:efficiency_chain}
\end{figure}

\begin{multicols}{2}
\subsection{Auxiliary Network: Isolated DC-DC Converter}

While the traction inverter operates directly off the 72V High-Voltage (HV) bus, the vehicle's critical telemetry modules, lighting arrays, cooling fans, and primary ECU mandate a stable 12V Low-Voltage (LV) supply. Tapping a subset of series cells is fundamentally forbidden as it guarantees fatal pack unbalancing. Instead, a dedicated DC-DC step-down converter must bridge the HV and LV domains.

\textbf{Topology Selection:} A High-Frequency Isolated Buck Converter (specifically a Forward Converter topology) is selected. Galvanic isolation via a high-frequency transformer is mandated by safety regulations to prevent catastrophic 72V potential from riding onto the exposed 12V driver chassis in the event of a semiconductor failure. 

\textbf{Working Explanation:} The 74V nominal DC input is chopped into a square wave via a power MOSFET switching at 100 kHz. This AC waveform steps down through the high-frequency ferrite transformer ($N_1:N_2$ ratio of approximately 5:1). The secondary 15V AC waveform is subsequently rectified by ultra-fast Schottky diodes and heavily smoothed via an LC low-pass filter to yield a pure 12V DC output. 

\textbf{Component Selection and Verified Calculations:}
Assume the LV system demands a peak 20 Amperes of current, requiring a 240W continuous DC-DC output. The exact required duty cycle ($D$) of the primary side MOSFET—assuming an idealized continuous conduction mode (CCM)—is fundamentally linked to the transformer turns ratio ($n = N_1/N_2 = 5$) and input potential ($V_{in} = 74\,\mathrm{V}$):

\begin{equation} \label{eq:duty_cycle}
\begin{split}
D &= \left( \frac{V_{out}}{V_{in}} \right) \cdot n \\
&= \left( \frac{12\,\mathrm{V}}{74\,\mathrm{V}} \right) \cdot 5 \\
&\approx 0.81 \ (81\%)
\end{split}
\end{equation}

For the critical output LC filter, the target peak-to-peak inductor current ripple ($\Delta I_L$) is specified as 20\% of the maximum 20A load ($\Delta I_L = 4\,\mathrm{A}$). The requisite output inductance ($L_{out}$) to suppress this ripple at a switching frequency ($f_{sw}$) of 100 kHz is verified as:

\begin{equation} \label{eq:inductor_calc}
\begin{split}
L_{out} &= \frac{V_{out} \left( 1 - \frac{D}{n} \right)}{\Delta I_L \cdot f_{sw}} \\
&= \frac{12 \cdot \left( 1 - \frac{0.81}{5} \right)}{4 \cdot 100,000} \\
&= \frac{12 \cdot 0.838}{400,000} \\
&\approx 25.1\,\mu\mathrm{H}
\end{split}
\end{equation}

Consequently, a 33 $\mu$H shielded power inductor is conservatively selected alongside a 100 V, 10 m$\Omega$ primary N-channel MOSFET, completing the highly efficient and isolated auxiliary power backbone of the vehicle.
\end{multicols}
